\chapter{绪论}{Introduction}

\section{本课题研究背景及研究意义}{
  Background and significance of research
}

随着电子商务、现代物流以及智能制造技术的快速发展，仓储系统在现代供应链中的作用日益突出。库存管理作为仓储作业中的基础环节，直接影响企业的出入库效率、库存准确性以及整体运营成本。自动化立体仓库（Automatic Storage and Retrieval System，AS/RS）因具有空间利用率高、作业效率高和自动化程度高等优势，已逐步广泛应用于烟草、医药、汽车零部件等多个行业。尤其在烟草仓储场景中，货物通常呈规则化堆叠存放，库存信息更新频繁，这对盘点工作的准确性、时效性和自动化水平提出了更高要求。

然而，传统库存盘点方法主要依赖人工完成，不仅劳动强度大、效率较低，而且在货物数量较多、堆叠结构复杂或作业环境受限的情况下，极易出现漏盘、错盘等问题。相较之下，精准、高效且自动化的库存盘点技术能够显著降低人工成本，提高盘点效率，并为后续仓储调度、库存预测和资源配置提供可靠的数据支撑。因此，研究适用于立体仓库场景的智能盘点方法，具有重要的现实意义和应用价值。

目前，仓储盘点领域主要存在以下几类技术路径。第一类方法基于 RFID（射频识别）技术。该方法能够实现快速批量识别，并具有较强的耐用性和较好的自动化水平。然而，其实施过程通常需要为货物配置电子标签，系统建设与后期维护成本较高，尤其在高吞吐量、大规模仓储场景中，这一成本问题更加突出，从而限制了其广泛推广应用。第二类方法主要采用 UAV（无人机）或 AGV（自动导引车）结合条码、二维码等识别码完成盘点作业。与 RFID 方法相比，此类方案在一定程度上降低了硬件成本，但仍需依赖外部标识介质，且在仓储现场运行过程中容易受到设备协同、视角变化和识别条件限制的影响，难以充分满足复杂场景下的自动化盘点需求。除上述两类方法外，部分仓储系统还尝试采用基于视觉或重量等传感器的辅助测量方式，例如通过比较货物在存储前后的图像或重量变化实现状态监测。当变化量超过预设阈值时，再由工作人员进行进一步确认和处理。该类方法在某些特定场景下具有一定实用性，但其对环境噪声、传感器精度以及现场工况变化较为敏感，系统稳定性有限，且仍然难以完全摆脱人工干预。

随着计算机视觉和深度学习技术的迅速发展，利用视觉感知方法实现库存盘点已成为智能仓储研究的重要方向。从计算机视觉角度来看，库存盘点任务可视为对仓库场景图像中的目标进行检测、分割以及数量统计的综合问题。相较于依赖标签或外部传感器的方法，视觉方法能够直接基于现场图像完成信息提取，具有更好的灵活性、较低的系统改造成本以及更强的场景适应性。对于立体仓库而言，现有监控系统已能够持续采集高分辨率图像，这为基于视觉的智能盘点提供了良好的数据基础。

在实际应用中，仓库管理系统通常要求盘点任务具备较高的实时性。如果将前端采集到的大量高清图像全部传输至远程云端进行处理，不仅会占用较大网络带宽，还可能影响物流系统中其他关键指令的正常传输。基于这一考虑，将图像处理任务部署在仓库现场的边缘服务器上是一种更具可行性的方案。边缘服务器能够在本地完成图像识别与数量估计，仅将最终识别结果传输至仓库管理系统，从而有效降低网络压力并提升系统响应效率。

然而，边缘部署环境中的计算资源通常较为有限，这使得高精度识别模型的大规模参数量与边缘侧轻量化部署需求之间形成明显矛盾。一方面，复杂堆叠场景中的遮挡、重叠、纹理相似和层次结构差异，使得普通检测模型容易出现漏检与误检，从而影响后续计数结果。另一方面，即便采用高精度深度学习模型，也往往面临参数规模大、计算量高、推理延迟较长等问题，难以直接满足边缘部署条件。此外，堆叠货物盘点任务并非单纯的目标识别问题，还涉及对上下层结构关系的理解以及基于结构规则的数量推理，这进一步提升了任务难度。

因此，如何在复杂堆叠场景下构建一种既具有较高检测与分割精度、又能够满足边缘部署轻量化需求，并进一步支持货物数量准确推断的智能盘点方法，成为当前立体仓库视觉盘点领域亟待解决的关键问题。围绕这一问题展开研究，不仅能够提升立体仓库盘点效率与库存管理水平，也能够为视觉感知在智能仓储中的落地应用提供新的理论支持和技术路径。

\section{国内外相关研究现状}{Analysis of the research status at home and abroad}

围绕立体仓库堆叠货物盘点任务，本文的研究内容主要涉及目标检测与实例分割、对象计数、模型压缩等方向。相关研究的发展为本文方法设计提供了重要参考，但现有方法在复杂堆叠场景适应性、结构化计数能力以及边缘部署效率等方面仍存在不足。下面从这几个方面对国内外相关研究现状进行分析。

\subsection{目标检测与实例分割}{}

堆叠货物数量统计的前提是对货物目标进行准确定位和轮廓提取，因此目标检测与实例分割是实现自动盘点的基础技术。目标检测作为计算机视觉领域的重要研究问题，主要包含目标定位与类别识别两个任务。实例分割则是在目标检测基础上的进一步扩展，不仅要求识别目标类别和位置，还需要输出每个实例的像素级掩膜，从而为后续结构分析与数量估计提供更精细的几何信息。

现有目标检测方法大致可分为两阶段方法、一阶段方法以及基于 Transformer 的方法。两阶段检测方法以 R-CNN、Fast R-CNN 和 Faster R-CNN 为代表，其核心思想是先生成候选区域，再完成分类与边界框回归。这类方法通常具有较高的检测精度，但由于处理流程较长，模型结构相对复杂，推理效率往往受到限制。一阶段目标检测方法如 SSD、RetinaNet 和 YOLO 系列，则通过单次前向传播直接完成目标定位与分类，具有结构简洁、推理速度快等优势，因而更适用于实时检测任务和边缘部署场景。

近年来，基于 Transformer 的目标检测方法受到了广泛关注。DETR 通过端到端学习框架避免了传统候选框生成和复杂后处理过程，但其在训练收敛速度和小目标检测性能方面仍存在不足。为解决这些问题，研究者提出了 Deformable DETR、Group DETR、DAB-DETR、DINO 等一系列改进方法，这些方法在训练效率和检测性能方面均取得了不同程度提升。但总体来看，基于 Transformer 的检测方法通常模型规模较大、训练成本较高，在边缘计算资源受限的场景中仍面临部署挑战。

对于实例分割任务，Mask R-CNN 是具有代表性的经典方法之一，其通过在检测分支之外增加掩膜分支，实现了较高质量的实例级分割。随后，YOLACT、SOLO、SOLOv2、TensorMask、QueryInst、SparseInst、Mask2Former 以及 YOLO 系列实例分割模型等方法，分别从单阶段推理、动态掩膜生成、稀疏表示和全局特征建模等角度推动了实例分割技术的发展。总体而言，单阶段实例分割方法在推理效率方面更具优势，而基于 Transformer 的分割方法在全局建模能力和掩膜表达能力方面表现更强，但其模型复杂度和部署成本也相对较高。

对于立体仓库堆叠货物盘点任务而言，方法不仅需要具备较好的目标定位与轮廓分割能力，还需要兼顾边缘部署条件下的推理效率与计算资源约束。因此，选择具有轻量化潜力的实例分割框架作为基线，并在此基础上开展针对复杂堆叠场景的结构改进，具有较强的现实必要性。

\subsection{对象计数}{}

对象计数是计算机视觉中的重要研究方向，常见方法大致可分为基于回归的方法、基于密度图的方法和基于检测的方法。基于回归的方法通常通过建立图像特征与目标数量之间的直接映射关系完成计数，其实现形式较为直接，但往往缺乏良好的可解释性。基于密度图的方法常用于高密度场景，其基本思路是先预测目标密度分布，再通过积分或累加得到总数量。该类方法在目标密集分布场景中具有一定优势，但对于目标的层次结构和个体边界表达相对不足。基于检测和分割的方法则通过识别单个目标实例实现数量统计，具有较强的可解释性，但其性能高度依赖于检测与分割结果的准确性。

立体仓库中的堆叠货物计数任务与一般对象计数任务存在明显差异。该任务不仅要求估计目标总数量，还涉及上下层结构关系判别以及基于堆叠规则的数量推理。对于此类场景，单纯依赖回归方法或密度图方法难以充分利用空间结构信息，而传统基于检测的方法在面对堆叠遮挡和层次关系时又存在明显不足。因此，面向堆叠货物盘点任务，仍有必要结合实例分割结果设计更具结构感知能力的后处理计数方法，以提升数量推断的准确性和可解释性。

\subsection{模型剪枝与模型蒸馏}{}

随着深度神经网络模型规模不断增加，模型压缩逐渐成为边缘智能领域的重要研究方向。其中，模型剪枝是最常见的压缩技术之一，主要包括非结构化剪枝和结构化剪枝两类。非结构化剪枝通常根据参数重要性删除部分连接或权重，可以显著降低参数冗余，但由于生成的稀疏结构不规则，实际硬件加速效果往往有限。结构化剪枝则通过删除整个通道、卷积核或模块来降低模型复杂度，其输出结构更加规整，更有利于在实际部署中获得真实的推理加速收益。

在结构化剪枝研究中，常见策略包括组级剪枝、滤波器剪枝和通道剪枝等。相关研究通常根据权重范数、激活值、能耗贡献或特征重要性等指标评估不同结构单元的重要程度，并据此执行剪枝操作。总体来看，结构化剪枝在实际部署友好性方面相较非结构化剪枝更具优势，因此更适合作为边缘侧模型压缩的重要手段。

知识蒸馏的核心思想是利用性能较强的教师模型指导轻量化学生模型训练，以实现“以大带小”的知识迁移。现有蒸馏方法包括基于输出响应的蒸馏、基于中间特征的蒸馏以及基于关系约束的蒸馏等。大量研究表明，知识蒸馏能够在一定程度上弥补轻量化模型在压缩过程中的性能下降问题，因此常与模型剪枝配合使用，以在模型复杂度和性能之间取得更优平衡。

综上，模型剪枝与知识蒸馏为深度模型在边缘部署场景中的压缩与加速提供了有效技术路径。其中，结构化剪枝更有利于实际部署过程中的推理加速，而知识蒸馏能够在一定程度上缓解压缩过程带来的性能退化。然而，针对堆叠货物盘点任务，如何在保持检测、分割与计数性能的前提下进一步降低模型复杂度，仍需开展针对性的研究。

\section{本文主要研究内容与论文结构}{Main research content and thesis structure}

针对立体仓库堆叠货物盘点任务中存在的识别精度不足、复杂场景计数困难以及边缘部署受限等问题，本文围绕“复杂堆叠场景下的高精度盘点”和“面向边缘部署的模型轻量化”两条主线开展研究。本文的主要研究内容如下：

\par\noindent
（1）面向复杂堆叠场景的实例分割模型改进与后处理计数方法研究。针对现有模型在堆叠货物识别过程中易出现漏检、误检以及对结构信息建模不足的问题，本文以 YOLO11-Seg 为基线，设计适用于堆叠货物场景的改进实例分割模型，并结合实例分割掩膜与空间几何关系构建后处理计数方法，以提升库存盘点的准确性与鲁棒性。

\par\noindent
（2）面向边缘部署的模型压缩方法研究。针对高精度识别模型参数量较大、计算开销较高而难以直接部署于边缘侧的问题，本文进一步研究模型压缩方法，通过结构化剪枝降低模型复杂度，并结合知识蒸馏缓解压缩过程中的性能损失，从而提升模型在实际仓储场景中的部署可行性。

研究内容之间的逻辑关系如\autoref{fig:fig1}所示。
\begin{figure}[H]
    \centering
    \includegraphics[width=0.9\textwidth]{fig1.pdf}
    \bicaption{研究内容逻辑关系}{Logical relationship of the research contents}
    \label{fig:fig1}
\end{figure}

本文的章节安排如下：第1章为绪论，主要介绍研究背景、国内外研究现状、研究内容与论文结构；第2章围绕堆叠货物检测、分割与计数任务，提出改进的实例分割模型及后处理计数方法；第3章针对边缘部署需求，研究模型压缩方法，包括结构化剪枝与知识蒸馏；第4章通过实验对所提方法的有效性进行验证，并从精度、复杂度与计数性能等方面进行分析；第5章对全文工作进行总结，并对后续研究方向进行展望。

